% Compiled with LuaLaTeX
\documentclass[tate, book, paper=a5]{jlreq}
\usepackage{luatexja-fontspec}
\usepackage{tikz}
\usepackage{eso-pic}
\usepackage{geometry}

% Margin settings
\geometry{margin=20mm}

% Subtle bottom-center geometric design using eso-pic and TikZ
\AddToShipoutPictureBG{
    \begin{tikzpicture}[remember picture, overlay]
        % Define a very subtle color
        \colorlet{accentcolor}{gray!35!white}
        
        % Position at the bottom center of the page (15mm from the bottom edge)
        \begin{scope}[shift={(current page.south)}, yshift=15mm]
            % Horizontal hairline
            \draw[line width=0.2pt, color=accentcolor] (-10mm, 0) -- (10mm, 0);
            
            % Minimal Pentagon (Pointing upwards, symbolizing forward movement/will)
            \fill[color=accentcolor] 
                (0, 1.5mm) -- (-1.2mm, 0.2mm) -- (-0.8mm, -1.2mm) -- (0.8mm, -1.2mm) -- (1.2mm, 0.2mm) -- cycle;
        \end{scope}
    \end{tikzpicture}
}

\begin{document}

\chapter*{走れメロス}

メロスは激怒した。必ず、かの邪智暴虐の王を除かなければならぬと決意した。メロスには政治がわからぬ。メロスは、村の牧人である。笛を吹き、羊と遊んで暮して来た。けれども邪悪に対しては、人一倍に敏感であった。

きょう未明メロスは村を出発し、野を越え山越え、十里はなれた此のシラクスの市にやって来た。メロスには父も、母も無い。女房も無い。十六の、内気な妹と二人暮しだ。この妹は、村の或る律気な一牧人を、近々、花婿として迎える事になっていた。結婚式も間近かなのである。メロスは、それゆえ、花嫁の衣裳やら祝宴の御馳走やらを買いに、はるばる市にやって来たのだ。先ず、その品々を買い集め、それから都の大路をぶらぶら歩いた。メロスには竹馬の友があった。セリヌンティウスである。今は此のシラクスの市で、石工をしている。その友を、これから訪ねてみるつもりなのだ。久しく逢わなかったのだから、訪ねて行くのが楽しみである。

\end{document}